%%%%%%%%%%%%%%%%%%%%%%%%%%%%%%%%%%%%%%%%%%%%%%%%%%%%%%%%%%%%%%%%%%%
%ElectroX_Devoir.tex : Devoir des cours de Physique électronique & Électronique et mesures expérimentales au Département de physique, de génie physique et d'optique - FSG - Université Laval. Les fichiers de chaque exercice sélectionné sont insérés à l'aide du module subfiles.
%------------------------------------------------------------------
%Créé par Claudine Nì. Allen
%Contributions par -
%Compilateur pdfLaTeX, distribution TeX Live 2023
%Dernière modification : 28 septembre 2023
%
%ToDo : voir le Template correspondant.
%%%%%%%%%%%%%%%%%%%%%%%%%%%%%%%%%%%%%%%%%%%%%%%%%%%%%%%%%%%%%%%%%%%%

\RequirePackage[l2tabu, orthodox]{nag} %%Check for obsolete commands
\documentclass[english,french,12pt]{article}
\input{../aPreambleInput.tex}
%----------------------------------------------------------------
%### Layout and other format configuration ###
%\the\parindent{} , \the\value{equation} %%and other '\the' commands to see the current parameter, counter,... value
\geometry{  %%for geometry package
%    showframe,  %%uncomment to see lines delimiting the dimension boxes
    letterpaper,
    margin=0.75in}
\fancypagestyle{Devoir}{
    \fancyhf{}
    \renewcommand{\headrulewidth}{0pt}
    \fancyfoot[C]{-\thepage-} %%bottom center page numbering
    \AtEndDocument{\thispagestyle{empty}}} %%removes the last page number
\pagestyle{Devoir}
%\interfootnotelinepenalty=10000 %%uncomment if long footnotes are incorrectly spreading on several pages
\widowpenalty=300 %%increase for less likely widow lines
\clubpenalty=300  %%increase for less likely orphan lines
\setlength{\parskip}{1em plus0.4em minus0.2em} %%changing space between paragraphs
\setlength\parindent{0em} %%uncomment to force indentation of all paragraphs or remove completely with 0em
%\setlist[enumerate]{wide=0pt,widest=99,leftmargin=\parindent,labelsep=*} %%Global changes of the [enumerate] list
\setcounter{equation}{0} %%ensure resetting of the equation counter
\captionsetup{  %%for figure and table captions
    font=footnotesize,
    labelfont=bf,
    labelsep=period,
    margin=5em}
    
%---------------------------------------------------------
\begin{document}
%---------------------------------------------------------
\renewcommand{\onlyinsubfile}[1]{}
%
%------ BEGIN HOMEMADE TITLE ------%
%inspiré de l'équipe du Prof. L.J. Dubé
%
\begin{center}
    \textbf{\large{Électronique et mesures expérimentales / Physique électronique}}\\
    \vspace{0.2em}
    \textbf{GPH-2006 / PHY-2002}\par
    \textsc{Département de Physique, de Génie Physique et d'Optique\\
    Faculté des Sciences et de Génie, Université Laval}
\end{center}
%\vspace{-1em}
\noindent Automne 2025 \hfill Responsable: C. Allen
%rajouter une * une *footnote qui link directement en renvoyant à la fin pour des instructions Gradecope que les co-équipiers se s'ajoutent après soumission (voir message Teams pour lien de vidéo, etc.) et que ce doit tjrs être la même personne pour éviter de générer 2 copies du même devoir quand on fait plusieurs soumissions.
\vspace{0.2em}
\hrule
%\vspace{-1.5ex}
%\centering 
\begin{center}
    \textsc{Devoir \# 2}\par
    \small{à remettre \textbf{en binôme}, \textit{i.e.} uniquement avec votre co-équipier.ère des laboratoires.}
\end{center}
%\vspace{-0.5ex}
%Début: 23 septembre 2024 \hfill Fin: 16 octobre 2024\par
%\vspace{0.4ex}
\hrule
\justify

%------ END HOMEMADE TITLE ------%
%
%\vspace{-0.5em}

\begin{enumerate}[label=\Roman*., wide, labelwidth=!, labelindent=0pt]
    \setcounter{enumi}{0}
    \vspace{2ex}
    \item \subfile{DevoirExa/AC transitoire ou permanent/AC_005}
    \newpage
    \item \subfile{DevoirExa/AC transitoire ou permanent/AC_004}
\end{enumerate}
 
 %------ BEGIN HOMEMADE FOOTER ------%
\vfill
\hrule
\vspace{0.3em}
\centering
\textbf{N'oubliez pas d'expliciter votre démarche et/ou de citer vos références.}\par
\vspace{-0.3em}
Sur \href{https://www.gradescope.com/}{Gradescope}, ajoutez votre coéquipier.ière après avoir sélectionné les pages correspondant à chaque numéro en soumettant le devoir une première fois.\par
\vspace{1em}
\hrule
\end{document}